\documentclass[../main.tex]{subfiles}
\begin{document}
%===================TIÊU ĐỀ TẬP=============================
\begin{table}[h]
  \begin{adjustwidth}{-1cm}{}
    \begin{tabular}{>{\centering\arraybackslash}p{6cm}|>{\raggedright\arraybackslash}p{14cm}}
      \multirow{5}{6cm}
      % Cột bên trái: ÉPISODE và số tập
      {\\[-40pt]
        \flaregothic\fontsize{20pt}{20pt}\selectfont \centering
        \color{eptype}PERSONNALISÉ\color{black}\\[12pt]
        \flaregothic\fontsize{36pt}{36pt}\selectfont \centering
        \color{epnum}{\vnmdisp KHÔNG}\color{black}\\[6pt]
        \cafeteria\fontsize{20pt}{20pt}\selectfont\color{epnum}(C-0)\color{black}
        \\[18pt]
        \flaregothic\fontsize{14pt}{14pt}\selectfont
        \color{epattr}BIENVENUE, \color{black}{\uline{Kuon} !}
        \color{black}}
       &
      % Dòng thứ nhất: tên tập tiếng Nhật
      {\fotskip\fontsize{18pt}{18pt}\selectfont \ruby{火}{ひ}\ruby{川}{かわ}の\ruby{始}{はじ}まり
      }                                                    \\[6pt]
      % Dòng thứ hai: tên tập tiếng Anh
       & {\cafeteria\fontsize{18pt}{18pt}\selectfont The Story Begins\dots? }                                                            \\[4pt]
      % Dòng thứ ba: tên tập tiếng Việt
       & {\vnmsans\fontsize{18pt}{18pt}\selectfont Ngày định mệnh của Hikawa Kuon}                                                       \\[8pt]
      % Dòng thứ tư: Nhân vật xuất hiện
       & {\caviar{\fontsize{14pt}{14pt}\selectfont Nhân vật: \emp{kuon1} \emp{achan} \emp{roboco} \emp{fubuki} \emp{matsuri} \emp{aqua}}} \\[8pt]
      %        }\\[8pt]
      %        % Dòng cuối cùng: Thời gian diễn ra của tập (thường sẽ là ngày phát hành)
      %        & {\continuum\fontsize{16pt}{16pt}\selectfont Ngày phát hành: 11/01/2019}
    \end{tabular}
  \end{adjustwidth}
\end{table}
%===========================================NỘI DUNG==========================================
\color{black}

\txtbx[2]{
  {\color{vol} \uline{Tóm tắt:}} Trong một buổi chiều tẻ nhạt như mọi ngày, Hikawa Kuon vô tình nộp đơn ứng tuyển vào vị trí nhân viên IT thử nghiệm của hololive production. Sau buổi phỏng vấn trực tuyến đầy tiếng cười và vô cùng hỗn loạn cùng CEO Tanigo, A-chan cùng bốn cô gái thế hệ đầu, cậu bất ngờ nhận được lời mời đảm nhận cùng lúc hai vai trò: kỹ thuật viên IT kiêm chức tổng quản lý của công ty. Dù chưa từng có kinh nghiệm quản lý thần tượng, Kuon vẫn quyết định dấn thân vào công việc đầy thử thách này để tự lo học phí đại học, chính thức bước vào một hành trình đầy tiếng cười và những tình huống oái oăm phía trước.
}

\begin{center}
  {\textit{*Các nhân vật nói bằng tiếng Etsunan.}}
\end{center}

Mùa xuân, năm 2019.

Tại quốc gia Etsunan - một thế giới song song kỳ lạ nơi tiếng Nhật được công nhận là ngôn ngữ chính thức thứ hai chỉ sau tiếng Etsunan bản địa, mọi thứ khác vẫn vận hành hệt như Việt Nam mà chúng ta đã biết. Nội thành \ruby{Kawachi}{Hà Nội} ở vũ trụ này có phần hiện đại hơn nhờ hệ thống tàu điện ngầm chạy êm ru xuyên lòng thành phố - thứ phương tiện công cộng vô cùng tiện lợi mà người dân có thể đi tới bất kỳ đâu, kể cả xuyên sang các thành phố lân cận khác.

Chiều hôm đó không có gì đặc biệt, giống như bao ngày bình thường khác trong căn nhà cao tầng nằm giữa khu đô thị mới ở phường \ruby{Anwa}{Yên Hòa} (từng thuộc quận \ruby{Kyuukami}{Cầu Giấy} cũ).

\dots Chính xác hơn, đó là một buổi chiều chẳng khác gì những ngày đã trôi qua trong suốt mười năm qua, kể từ khi Hikawa Kuon (\ruby{火}{ひ}\ruby{川}{かわ}\ruby{久}{く}\ruby{遠}{おん}, Hỏa Xuyên Cửu Viễn) bắt đầu duy trì thói quen tẻ nhạt của mình. Cậu nằm ườn trên tấm đệm, buông một tiếng thở dài thườn thượt, đôi mắt vô hồn nhìn lên trần nhà: ``Lại một ngày nữa trôi qua chẳng khác gì hôm qua. Ngủ, ăn, chơi game, rồi lại ngủ\dots\ Cái vòng lặp vô tận này bao giờ mới chịu kết thúc đây? Cứ thế này mãi thì mình mốc meo mất.''

Sở dĩ cậu rơi vào hoàn cảnh này là vì đang tận hưởng một năm nghỉ ngơi (\textit{gap year}) sau kỳ thi đầu vào đại học đầy áp lực và mệt mỏi, đặc biệt là khi cậu vừa mới xuất viện không lâu. Tuy rằng nó không căng thẳng tột độ như hồi thi vào trung học phổ thông, nhưng sức nóng của kỳ thi đại học vừa rồi cũng đã vắt kiệt sức lực của cậu.

Cuộc sống thường ngày của cậu cứ xoay quanh những hoạt động đơn điệu, lặp đi lặp lại không chút biến đổi. Bộ quần áo sơ-mi nhăn nhúm phối cùng chiếc quần thể dục xập xệ, cái bụng hơi phệ nhẹ cùng mái đầu rối bù xù như tổ quạ dường như là quá đủ để khắc họa bức chân dung lười nhác của chàng trai hai mươi hai tuổi này. Độ tuổi đáng ra phải ngập tràn hoài bão, ước mơ và hy vọng, nhưng với cậu thì chỉ có sự lười biếng và cái danh ăn bám bố mẹ. Gọi là ``thảm hại'' thì hơi quá, nhưng cũng có phần đúng, vì ngoài tài vọc vạch máy tính ra thì cậu thật sự rất, rất lười.

Tuy vậy, nhịp sống đơn điệu đó dường như đã sắp đi đến hồi kết. Đến tháng Năm này, cậu sẽ bắt đầu nhập học vào ngôi trường đại học đã chọn - Trường Đại học Xây dựng Kawachi, theo đuổi ngành Công nghệ thông tin. Chỉ cần nhìn cái cách cậu lúc nào cũng táy máy với đống linh kiện điện tử hay say mê phá đảo những trò chơi quen thuộc trên máy tính là đủ biết cậu yêu công nghệ đến nhường nào.

Gia đình Hikawa chung sống trong căn nhà sáu tầng chỉ rộng chưa đầy sáu mươi mét vuông, bao gồm bố mẹ cậu, ông bà nội, người bác là anh trai thứ của bố cậu, và cô em gái kém cậu mười tuổi.

Nằm trên tấm đệm với một tay đang lướt báo trên điện thoại, một tay gác lên trán, ánh mắt lơ đãng của cậu dõi qua từng dòng tin tức. Nhịp sống đều đều tẻ nhạt bắt đầu khiến cậu cảm thấy ngột ngạt: ``\textit{Mình thật sự nên kiếm một công việc làm thêm\dots\ Mình cũng sắp học đại học rồi, không thể cứ ngửa tay xin tiền bố mẹ với em gái mãi được.}'' Có vẻ như cậu cũng đã bắt đầu ý thức được việc tự lo cho bản thân.

``\textit{Nhưng vấn đề là mình biết làm cái gì bây giờ?}'' Ý nghĩ tự lập vừa nhen nhóm lập tức vấp phải sự lười biếng vốn có. ``\textit{Không biết liệu mình có phải mãi chật vật thế này không nữa\dots}'' cậu lẩm bẩm đầy ngán ngẩm.

Cậu tiếp tục lướt ngón tay trên màn hình điện thoại, rồi đôi mắt vô tình dừng lại trước một mẩu quảng cáo lớn chiếm gần như toàn bộ diện tích hiển thị.

Ban đầu, theo phản xạ tự nhiên của một cư dân mạng, cậu định vuốt tay đóng nó đi vì nghĩ rằng ``nó cũng chỉ giống như những quảng cáo vặt vãnh, rác rưởi không đáng bận tâm''. Thế nhưng, hình ảnh phối màu rực rỡ cùng dòng tiêu đề bắt mắt đã khiến ngón tay cậu chững lại giữa chừng.

Một chiếc logo gồm hai hình tam giác màu xanh đậm và nhạt lồng vào nhau, với dòng chữ ``hololive production'' xuất hiện đầy nổi bật ở phía dưới. Chiếc logo đơn giản nhưng mang lại cảm giác cực kỳ chuyên nghiệp, kèm theo dòng chữ tuyển dụng đầy khiêu khích: ``Tuyển dụng nhân viên IT - Cơ hội đồng hành cùng một trong những dự án đổi mới giải trí hàng đầu Nhật Bản!''

Kuon nhíu mày, sự hứng thú vốn đã nguội lạnh từ lâu đột nhiên trỗi dậy mạnh mẽ: ``Tuyển dụng vị trí IT à? Nghe cũng có vẻ ra gì và này nọ đấy chứ\dots''

Dòng chữ nhỏ hơn bên dưới tiếp tục hé lộ: ``Gia nhập đội ngũ sáng tạo của hololive - nơi hiện thực hóa giấc mơ giải trí ảo.''

``Cơ mà, đây là công ty giải trí ư? Sao nghe có vẻ giống như Love Live vậy nhỉ?'' cậu tự vấn đầy hoài nghi. Gạt bỏ ý nghĩ đó sang một bên, cậu bấm vào liên kết trang web quảng cáo, tự nhủ: ``Thôi kệ, cứ thử vào xem thế nào, biết đâu lại là một trang web lừa đảo công nghệ cao.''

Nhìn kỹ một lượt, giao diện trang web trông vẫn còn khá sơ khai và đơn giản, nhưng sau khi Kuon cẩn thận kiểm tra tên miền thì thấy nó cực kỳ uy tín và rõ ràng. ``Trông vẫn hơi nghi nghi\dots\ Hy vọng mình không bị ăn một quả lừa ngoạn mục,'' cậu vừa lẩm bẩm và nhấp chuột.

Cậu bấm vào biểu tượng menu ở góc màn hình, một thanh menu phụ hiện ra với đầy đủ các mục thông tin được xây dựng chỉn chu. ``Ghê thật, trang này còn thiết kế cả phần giới thiệu chi tiết nữa chứ\dots\ Để xem họ viết cái gì nào\dots''

Những dòng chữ liền hiện ra, từng dòng, từng dòng một. Đây rõ ràng là một trong những trang web được thiết kế theo phong cách hiện đại, trực quan dù quy mô công ty lúc này có vẻ vẫn còn rất nhỏ. ``Ờm\dots\ là một công ty con của Cover Corp., chuyên đào tạo những thần tượng ảo (Virtual Idol)\dots\ Hmm, thú vị đấy.''

Dòng chữ và những hình ảnh đồ họa đầy cuốn hút khiến cậu không thể dừng việc kéo chuột xuống dưới.

``Xem nào\dots\ Còn phân ra cả thế hệ nữa này. Thế hệ 0, thế hệ 1, thế hệ 2, GAMERS, INoNaKa\footnote{Một nhánh phụ của Cover Corp., cùng với hololive, bao gồm hai thành viên là Hoshimachi Suisei (sẽ được giới thiệu ở quyển số Bốn) và AZKi (sẽ được giới thiệu ở quyển Mười Một). Sau này họ chuyển về Gen 0 của \textit{hololive} khi nhánh bị giải thể vào giữa năm 2020.}\dots\ Sao hai cái cuối nghe tên trông lạc quẻ thế nhỉ?''

Tò mò trỗi dậy, cậu quyết định dành thêm thời gian tìm hiểu sâu hơn. Sau một hồi tra cứu thông tin trên các diễnàm mạng, cậu đã nắm được một vài nét cơ bản về công ty giải trí kỳ lạ này. Tokino Sora - thần tượng ảo đầu tiên của họ - chính là người đã đặt viên gạch đầu tiên cho dự án này với một buổi phát sóng trực tiếp ban đầu chỉ có vỏn vẹn 13 người xem trực tuyến. Câu chuyện đầy tính truyền kỳ này nhanh chóng thu hút sự chú ý của cậu.

``Hội nghị bàn tròn 13 người\dots\ Nghe cũng hoành tráng và đầy cảm hứng đấy chứ.''

Quay trở lại trang tuyển dụng, trước mắt cậu bây giờ là một biểu mẫu đăng ký kèm theo các trường thông tin cần điền. Trên đó, họ yêu cầu ứng viên điền đầy đủ họ tên, quê quán, email và thậm chí là cả số căn cước công dân. Đến lúc này, bản năng phòng thủ của một dân IT bắt đầu rung chuông cảnh báo.

``Không lẽ là lừa đảo thu thập thông tin cá nhân? Mình chưa từng thấy có trang web tuyển dụng chính thống nào lại yêu cầu điền cả số căn cước ngay từ bước nộp đơn thế này!''

Tuy nhiên, sau khi kiểm tra lại giao thức bảo mật và xác minh lại liên kết chính chủ của Cover Corp, cậu nhận ra đây thực sự là trang web chính thức của công ty chứ không phải trang giả mạo. ``Mà thôi kệ đi, cùng lắm chắc là họ sẽ gọi điện làm phiền hoặc khủng bố tin nhắn rác thôi chứ gì. Cái đó thì mình chặn số phút mốt,'' cậu tặc lưỡi tự trấn an.

Quyết tâm dẹp bỏ những nghi ngờ cuối cùng trong lòng, Kuon lấy hết can đảm điền tất cả thông tin cá nhân vào biểu mẫu rồi dứt khoát nhấn nút gửi đơn.

Màn hình lập tức hiển thị thông báo phản hồi: ``Đăng ký thành công. Xin hãy chờ trong vài giờ để hệ thống xử lý. Một email xác nhận sẽ được gửi đến tài khoản của bạn khi hoàn thành.'' Cậu thầm nghĩ đầy lo lắng: ``\textit{Xử lý nhanh thế! Nhưng mà\dots\ trông có vẻ uy tín\dots\ chắc vậy?}''

``\textit{Nhưng mà giờ mình nên làm gì tiếp theo đây? Đành phải để số phận đưa đẩy thôi chứ biết sao giờ\dots\ Để xem cái công ty giải trí này còn cái gì hay ho nữa không nào\dots}''

Kuon tiếp tục lướt web để tìm hiểu thêm về cái gọi là ``hololive'' này. Cậu phát hiện ra khá nhiều thông tin bên lề thú vị, từ những khoảnh khắc hài hước (clipping), các thước phim hoạt hình ngắn của các thành viên, cho đến cả những bí mật ít người biết đến.

``Thế hệ 1 từng có một thành viên khác nữa à? Để xem\dots\ Hitomi\dots\ Chris?'' cậu tò mò lẩm bẩm đọc to thông tin trên màn hình. ``Ờm\dots\ bị sa thải vì vi phạm hợp đồng nghiêm trọng chỉ ngay sau vài ngày debut\dots\ ồ.'' Cậu tiếp tục rê chuột tìm kiếm thông tin liên quan, rồi đôi mắt chợt dừng lại trước những dòng bình luận của cộng đồng mạng.

``Bị dính lùm xùm tình cảm cá nhân\dots\ xong rồi làm rò rỉ cả thông tin bảo mật của công ty\dots\ À, chả trách,'' cậu lẩm bẩm rồi chẹp miệng lắc đầu, ``bị sa thải là đúng rồi. Bảo sao người ta gọi cô nàng này là một truyền thuyết đô thị của giới ảo.''

Cậu tiếp tục tra cứu thêm một lúc nữa cho đến khi đã nắm bắt được những khái niệm cơ bản nhất về công ty giải trí này, rồi mới chịu quay lại với thú vui quen thuộc của mình là đọc báo mạng để giết thời gian.

Thế nhưng, cậu chàng không hề hay biết rằng\dots

\dots quyết định điền đơn đầy ngẫu hứng trong một buổi chiều tẻ nhạt này sẽ chính là chiếc chìa khóa mở ra một chương mới vô cùng hỗn loạn, rực rỡ và không tưởng trong cuộc đời mình.

\ct

Buổi tối ngày hôm đó. Trong gian bếp nhỏ ấm cúng, hai bố con nhà Hikawa đang cùng nhau thưởng thức bữa cơm tối đơn giản. Mẹ và cô em gái út đã về quê chơi vài hôm nhân dịp nghỉ học kỳ, để lại căn nhà rộng rãi cho hai người đàn ông tự quản. Như thường lệ, hai bố con vừa ăn vừa rôm rả trò chuyện đủ thứ trên đời, từ chuyện học hành của cậu cho đến những vấn đề xã hội xa xôi, thân thiết hệt như hai người bạn già lâu năm.

Bất chợt, tiếng chuông thông báo điện thoại của Kuon vang lên giòn giã. Cậu liếc nhanh qua màn hình rồi reo nhẹ: ``Chắc là email phản hồi của người ta đến rồi.''

``Email gì thế con?'' Ryujin - bố của cậu - nhấp một ngụm nước, tò mò hỏi.

``Đơn xin việc của con ạ,'' cậu hào hứng mở khóa điện thoại, nhanh chóng kiểm tra thông báo mới nhất trên màn hình, ``À, họ duyệt đơn của con rồi này!''

Ông ngạc nhiên nhướng mày: ``Xin được việc rồi á? Sao nhanh thế con?''

``Thì mấy đứa tầm tuổi con bây giờ cũng bắt đầu đi làm kiếm thêm tiền tiêu vặt rồi bố,'' cậu tặc lưỡi đáp lại.

``Tốt quá! Thế con ứng tuyển vào vị trí gì đấy?''

``IT ạ,'' cậu vừa liếc qua thông tin đăng ký vừa giơ màn hình điện thoại về phía bố, ``Kiểu mấy công việc đụng đến máy tính, mạng mẽo ấy mà.''

``Có yêu cầu gì cao không?''

``Yêu cầu lập trình cơ bản thôi ạ, với cả có hiểu biết tốt về máy móc và thiết bị công nghệ thông tin. Nói chung là đúng chuyên môn vọc vạch của con rồi.''

``Đúng sở trường thế là tốt. Vậy lương bổng người ta trả thế nào?''

``Họ ghi là đang cần người gấp, nên con đoán chắc không dưới mười nghìn yên một tháng đâu bố.''

Ryujin nheo mắt nhìn con trai: ``Chắc gì. Nhỡ tháng này họ cần tuyển gấp nhưng tháng sau hết việc thì sao?''

``Đấy là con đoán thế thôi\dots\ À mà bố ơi, họ gửi tiếp một email hướng dẫn rồi này.''

``Họ nói sao con?'' ông đặt cốc bia xuống bàn, chăm chú nhìn cậu.

``Bảo là tầm chín giờ tối nay sẽ có một buổi gặp mặt và phỏng vấn trực tiếp qua Zoom. Họ gửi sẵn cả mã phòng họp đây rồi ạ.''

Ryujin nghe vậy thì trong lòng không giấu nổi niềm vui và sự tự hào khi thấy cậu con trai cả cuối cùng cũng bắt đầu có ý thức tự lập. Tuy thế, với kinh nghiệm của một người đi trước, ông vẫn cẩn thận dặn dò: ``Cẩn thận đấy con ạ, thời buổi này lừa đảo mạng nhiều lắm.''

``Con cũng đã nghi nghi từ chiều rồi,'' cậu nhún vai đáp, ``Nhưng mà chắc không sao đâu bố. Trang web con đăng ký là trang chính chủ của tập đoàn bên Nhật mà, con kiểm tra kỹ tên miền rồi.''

``Dù vậy thì con vẫn cứ phải giữ cái đầu lạnh và quan sát cho cẩn thận đấy.''

Cậu con trai kéo dài giọng đáp một tiếng ``Rồi~{}'' đầy lém lỉnh, trước khi đặt điện thoại xuống bàn để tiếp tục bữa cơm tối ấm cúng cùng bố.

\ct

Ăn xong bữa cơm, hai bố con cùng nhau dọn dẹp sạch sẽ gian bếp. Kuon nhanh chóng trở về phòng riêng, mở chiếc laptop quen thuộc và truy cập vào ứng dụng họp trực tuyến Zoom. Cậu cẩn thận nhập mã phòng như hướng dẫn trong email, chỉnh lại góc camera rồi bắt đầu ngồi chờ. Lồng ngực cậu khẽ phập phồng, không giấu nổi sự hồi hộp xen lẫn chút phấn khích.

``\textit{Không biết họ sẽ phân công việc cho mình kiểu gì nhỉ\dots}'' cậu thầm nhủ, đôi mắt không rời khỏi màn hình trống trơn.

Tuy nhiên, thời gian cứ thế trôi qua. Đã quá giờ hẹn khoảng mười phút mà phòng họp vẫn im lìm không một bóng người. Cậu bắt đầu nảy sinh nghi ngờ, nhíu mày tự hỏi: ``Không lẽ đây là một quả lừa đảo thật sao?'' Nhưng dù sao đi nữa, Kuon vẫn quyết định kiên nhẫn ngồi chờ thêm, tranh thủ mở một trò chơi điện tử lên để giết thời gian.

Ba mươi phút đằng đẵng trôi qua. Phòng họp trực tuyến vẫn cô quạnh chỉ có mình cậu ngốc nghếch ngồi trước màn hình. Cảm giác sốt ruột và lo lắng bắt đầu dâng lên rõ rệt.

Từ ngoài phòng khách, tiếng bố Ryujin vọng vào: ``Vẫn chưa vào được hả con?''

``Vào được từ nãy rồi bố ạ, nhưng ngoài con ra thì chẳng thấy có ai cả,'' Kuon nói vọng ra, giọng đượm vẻ thất vọng.

``Hay thử gọi điện thoại trực tiếp cho người ta xem thế nào?'' bố cậu hiến kế.

Cậu thở dài thườn thượt, lắc đầu dù biết bố không nhìn thấy: ``Khổ nỗi người ta còn không thèm ghi số điện thoại liên hệ trên website mới đau chứ bố. Kiểu này chắc con ăn quả lừa thật rồi--''

Chưa kịp để cậu dứt câu than thở, loa máy tính bỗng phát ra những tiếng rè rè chập kết nối, rồi một giọng nói cất lên từ phía bên kia màn hình: ``A-lô\dots\ A-lô? Cậu có nghe tôi nói không?''

Cùng lúc đó, giao diện Zoom liên tục hiện thông báo người dùng mới tham gia. Chỉ trong chớp mắt, phòng họp đã có thêm sáu thành viên mới. Năm nữ, một nam.

\begin{center}
  {\textit{*Các nhân vật nói chuyện với nhau bằng tiếng Nhật.}}
\end{center}

Trong số năm người phụ nữ xuất hiện trên camera, phần lớn trông đều còn rất trẻ, thậm chí có người cảm giác còn ít tuổi hơn cả Kuon. Người đàn ông duy nhất trong số họ thì có vẻ ngoài chững lại, trông đã ngoài ba mươi tuổi. Dựa vào những thông tin tìm hiểu được từ chiều, Kuon lập tức nhận ra bốn người trong số họ. Người phụ nữ còn lại có mái tóc màu xanh dậm dài ngang cổ, đằng sau buộc gọn bằng một chiếc nơ màu xanh lam xinh xắn. Cô đeo một chiếc kính gọng đen, đôi mắt xanh nước biển thông minh ẩn sau tròng kính, mặc chiếc áo phông đen đơn giản.

``Có, có ạ! Em nghe rất rõ,'' Kuon vội vàng điều chỉnh lại micro, cất tiếng trả lời bằng tiếng Nhật.

``Được rồi. Xin lỗi cậu nhiều nha, tại mạng bên văn phòng chúng tôi bất ngờ bị sập\dots'' giọng người phụ nữ đầy ái ngại.

Từ một khung camera khác, một giọng nói nhí nhảnh vang lên đầy vẻ hiếu kỳ: ``Ồ, một chàng trai trẻ à? Hiếm có khó tìm nha~''

Màn hình hiển thị của sáu người phía bên kia vẫn còn giật cục dữ dội do đường truyền mạng quốc tế kém ổn định. Khung hình của vị CEO Tanigo thi thoảng lại đứng hình mất vài giây với nụ cười hiền hậu thương hiệu bị đóng băng một cách đầy hài hước, trong khi âm thanh phát ra thì méo mó, gián đoạn. Phải mất chừng năm phút loay hoay căn chỉnh thiết bị, tín hiệu đường truyền mới tạm thời ổn định để họ có thể bắt đầu buổi phỏng vấn.

``Cậu là Hikawa Kuon, đúng chứ?'' người phụ nữ đeo kính cất tiếng hỏi để xác nhận thông tin.

Cậu vội vã gật đầu: ``Dạ đúng ạ. Còn chị là\dots''

``Tôi là Yuujin A, nhưng cậu cứ gọi tôi là A-chan cho thân mật nhé,'' cô nàng mỉm cười thân thiện qua màn hình.

``Dạ, em không dám đâu ạ. Chị là tiền bối của em, nên xin phép cho em được gọi là A-san nhé,'' Kuon lễ phép trả lời.

Cô nàng khẽ bật cười, nét mặt giãn ra: ``Thôi được rồi, cậu thích gọi thế nào cũng được. Tôi là người phụ trách quản lý nhân sự kiêm vận hành ở đây. Còn người đàn ông bên cạnh tôi đây chính là\dots''

Chưa đợi A-chan dứt câu giới thiệu, Kuon đã nhanh nhảu cướp lời với một nụ cười đầy ẩn ý: ``Cái này thì em biết rõ ạ. Là Motoaki Tanigo-san, hay còn được người hâm mộ và các thành viên gọi yêu thương là YAGOO - do bị Subaru-chan gọi nhầm tên trên sóng livestream đúng không ạ? CEO đáng kính của cả tập đoàn Cover.''

Nghe thấy thế, vị CEO trung niên lập tức bật cười xòa đầy bất ngờ và thích thú trước sự hiểu biết của cậu ứng viên trẻ. Trong khi đó, bốn cô gái còn lại ở các khung camera bên cạnh bắt đầu cười nghiêng ngả không ngừng, tạo nên một bầu không khí hỗn loạn vô cùng vui nhộn.

``Cậu nhóc này nói chuẩn không cần chỉnh luôn kìa! Dám chắc là cậu đã nghiên cứu cực kỳ kỹ lưỡng về công ty chúng tôi từ trước rồi đúng không?'' Tanigo hào hứng hỏi, mắt sáng lên đầy kỳ vọng.

``Trời đất ơi, Subaru-chan mà nghe thấy đoạn này chắc chỉ có nước khóc thét vì bị 'bóc phốt' xuyên quốc gia mất thôi!'' Aqua khẽ thốt lên. Giọng cô nàng có chút ngập ngừng e thẹn, đôi má hơi ửng hồng. Nhưng ngay lập tức, camera của Aqua bị giật lag, đóng băng đúng lúc cô đang dùng hai tay che mặt để lộ biểu cảm xấu hổ vô cùng đáng yêu.

``Chuẩn quá luôn rồi á!'' Fubuki gật gù tán thưởng, chiếc tai cáo trên đầu khẽ vẫy vẫy đầy phấn khích.

A-chan khẽ ho khan một tiếng để giữ lại chút nghiêm túc cuối cùng cho buổi phỏng vấn: ``Thôi nào mấy đứa, đừng có làm loạn lên thế chứ. Phải để lại ấn tượng tốt đẹp cho nhân viên mới chứ.''

``Không sao đâu, A-san,'' Kuon cười xòa, gãi đầu đầy thông cảm, ``Nhìn bầu không khí thế này là em cũng đủ hình dung ra ngày thường ở văn phòng náo nhiệt đến mức nào rồi.''

``Anh nói đúng chuẩn luôn ấy!'' Matsuri giơ ngón tay cái tán thành đầy hào hứng. Khung hình của cô nàng do giật lag nên đột ngột nhảy vọt lên với nụ cười toe toét cực kỳ nghịch ngợm sát sạt vào ống kính camera.

``Ủa, vậy nếu anh đã biết rõ về CEO của tụi em rồi, thì chắc anh cũng nhận ra bốn tụi em là ai rồi chứ nhỉ?'' Roboco nghiêng đầu hỏi, đôi mắt vàng của cô chớp chớp đầy tò mò.

Cậu gật đầu cười đáp: ``Tất nhiên rồi. Mà thực ra, kể cả anh không tra cứu từ trước thì tên hiển thị của mọi người trên Zoom vẫn sờ sờ ra đó thôi.''

Bốn cô gái đồng loạt nhìn xuống góc màn hình của mình, phát hiện ra tên thật của họ vẫn đang hiển thị lù lù bằng tiếng Nhật. Cả nhóm ngớ người ra một giây rồi đồng thanh cười trừ đầy ngượng ngùng.

``Cậu nhóc này cũng tinh mắt và nhạy bén ra phết đấy,'' YAGOO gật đầu gù tán thưởng.

Kuon khẽ thở dài nhẹ nhõm: ``Chuyện cơ bản mà chú Tanigo.'' Rồi cậu lập tức điều chỉnh lại nét mặt, nghiêm túc chuyển sang chủ đề chính của buổi họp: ``Mà thôi, đùa vui thế đủ rồi ạ. Vậy cụ thể thì công việc sắp tới của cháu ở công ty sẽ là gì ạ?''

``Cậu nhóc này nghiêm túc ghê ta. Thôi được rồi, để tôi giới thiệu lại một cách bài bản về công việc cậu sắp đảm nhận nhé\dots'' Tanigo hắng giọng, chính thức bắt đầu vào phần việc chính. ``Chắc cậu cũng đã biết, \hll\ mới được thành lập được hơn một năm, kể từ thời điểm Sora-chan ra mắt và đặt những viên gạch đầu tiên.''

``Dạ, cái đó cháu có tìm hiểu qua ạ.''

``Với ước mơ cháy bỏng của tôi là muốn thành lập một dàn thần tượng ảo vĩ đại và chuyên nghiệp giống như nhóm AKB48 phiên bản 2D, tôi đã và đang không ngừng nuôi dưỡng các tài năng ở nơi này\dots'' Tanigo say sưa thuyết trình về lý tưởng của mình.

Nghe đến đây, Kuon nhìn bốn cô gái đang chí chóe trêu chọc nhau trên màn hình Zoom, tới độ đổ mồ hôi hột: ``\textit{Mình thật sự không biết là với bốn cô nàng nghịch ngợm, tấu hài cực mạnh thế kia thì chú ấy đang nuôi dưỡng thần tượng hay đang đào tạo những chú hề mua vui nữa\dots\ Mà thôi, nghĩ trong đầu thế thôi chứ nói ra thì dở lắm.}''

``\dots Tính đến nay, chúng tôi đã phát triển được bốn thế hệ, trải dài từ thế hệ 0 đến thế hệ 2, và cả thế hệ GAMERS nữa,'' Tanigo tiếp tục giải thích.

``Cháu đoán thế hệ GAMERS chính là thế hệ 'hai phẩy năm' đúng không ạ?'' cậu hỏi.

A-chan gật đầu đồng ý: ``Đúng vậy đấy. Tôi nói thật lòng nhé, hồi đầu tôi cũng không dám tin là ông ấy có thể chèo lái và nuôi dưỡng được nhiều nhân sự cá tính đến như thế đâu.''

``Chắc hẳn là Tanigo-san và A-san đã phải vất vả và tốn nhiều tâm huyết lắm nhỉ,'' Kuon nói đầy cảm thông. Cả cô nàng Đeo kính và vị CEO trung niên đều gật gù đồng cảm sâu sắc: ``Đúng ha\dots'' / ``Không cãi vào đâu được\dots''

``Phải công nhận là các cô gái của chúng tôi đôi lúc cũng rất hay bày ra những trò oái oăm, nghịch ngợm vô cùng 'bá đạo' như vậy\dots'' Tanigo cười khổ đầy bất lực.

``Nhưng tôi nghĩ chính sự đa dạng và độc đáo đó mới là nét hấp dẫn riêng biệt tạo nên linh hồn cho công ty này,'' A-chan mỉm cười tiếp lời.

Nghe thấy thế, Matsuri lập tức phồng má, tỏ vẻ oan ức đầy nhí nhảnh: ``Ý chị là sao cơ, A-chan? 'Bá đạo' á? Em lúc nào cũng ngoan hiền, thuần khiết và là biểu tượng idol gương mẫu cơ mà!''

A-chan lập tức liếc xéo sang khung camera của Matsuri với vẻ mặt cạn lời: ``Thật không đấy? Thế lần trước ai đã bày trò trêu chọc giấu đồ của Nakiri-san suốt cả một tuần liền khiến em ấy suýt phát khóc hả?''

Matsuri cười gượng gạo, giơ hai tay lên phân bua đầy chống chế: ``Chị hiểu nhầm em rồi! Đấy là em đang giúp chị ấy rèn luyện thần kinh thép để chuẩn bị cho các buổi livestream lớn thôi mà!''

Kuon chứng kiến cảnh tượng đó thì chỉ biết cười thầm, lẩm bẩm: ``\textit{Xem ra hololive đúng là nơi để mọi người khai phá tiềm năng tấu hài thiên bẩm của mình rồi.}''

``Cậu không cần phải nói huỵch toẹt cái sự thật phũ phàng đó ra thế đâu\dots'' cô nàng Đeo kính đỡ trán thở dài.

CEO Tanigo mỉm cười hiền từ: ``Tôi nghĩ chính sự độc bản của mỗi thành viên đã tạo nên bầu không khí gia đình đặc biệt ở đây. Mỗi người một vẻ, không ai giống ai.''

A-chan, Kuon và Tanigo cùng nhìn bốn cô gái nghịch ngợm đang trêu đùa nhau qua màn hình và mỉm cười. Cậu con trai hắng giọng, gương mặt trở nên nghiêm túc và chuyên nghiệp hơn hẳn: ``Nhân tiện, bốn người có thể tự giới thiệu bản thân một chút được không ạ? Tuy em đã đọc qua hồ sơ của mọi người trên trang web rồi, nhưng em vẫn muốn được nghe trực tiếp từ mọi người để dễ làm quen hơn.''

Bốn cô gái gật đầu đồng ý và bắt đầu phần tự giới thiệu đặc trưng của mình.

``Wasshoi! Biểu tượng của sự thuần khiết vô ngần và là thần tượng quốc dân của mọi nhà, Natsuiro Matsuri đây!'' Matsuri nháy mắt tinh nghịch, làm điệu bộ vô cùng nhí nhảnh trước camera.

``Kon-kon-kitsune! Chúc anh một ngày tốt lành nhé! Em là Shirakami Fubuki đây ạ!'' Fubuki mỉm cười tươi tắn, giơ hai ngón tay tạo hình tai cáo cực kỳ chuyên nghiệp.

``Kon-aqua\~{} Em là Minato Aqua\dots\ Cơ mà anh đừng có nhìn em chăm chú qua camera thế chứ, xấu hổ chết đi được\dots'' Aqua cúi gầm mặt xuống bàn, hai tai đỏ ửng, giọng nói lí nhí đầy e thẹn hệt như một chú mèo con nhút nhát.

``HaroboHarobo\~{} Em là Roboco-san!'' Roboco vẫy tay chào đầy thân thiện.

Nhưng ngay lập tức, Matsuri đã chen ngang với nụ cười đầy châm chọc: ``Và cũng là 'đại PON' hậu hậu đệ nhất công ty của tụi em nữa đó!''

``Chị không có PON nha!!!'' Roboco phồng má giận dỗi. Do đường truyền mạng đột ngột bị bóp băng thông, tiếng hét phản đối của cô nàng qua bộ lọc âm thanh Zoom nghe giật cục và rè rè hệt như tiếng robot kim loại thực thụ: ``Ch-Chị\dots\ kh-không\dots\ P-PON\dots\ nhaaanhaaa\~{}\~{}!'' làm cả nhóm cười ồ lên.

Kuon không nhịn được cười, trêu chọc thêm: ``Đúng quá rồi còn gì nữa Roboco-san ơi. Mấy pha tự hủy đi vào lòng đất của em đầy rẫy trên mạng kìa, anh xem clip cắt mà cười đau cả bụng luôn đấy.''

Roboco phồng má giận dỗi, quay mặt đi chỗ khác tránh ống kính: ``Cả anh cũng hùa theo trêu em nữa sao!? Thật là bất công quá đi mà\~{}!''

Vị CEO của \hll\ mỉm cười, tiếp tục dẫn dắt câu chuyện: ``Chắc cậu cũng đã phần nào hiểu được tính cách của họ rồi nhỉ? Chúng tôi vẫn còn rất nhiều tài năng tuyệt vời khác nữa đấy!''

``Tiếc là hôm nay phần lớn các thành viên khác đều đang bận chuẩn bị cho buổi phát sóng trực tiếp của họ, nên chỉ có bốn đứa này là rảnh rỗi thôi,'' A-chan giải thích thêm.

Kuon gật đầu thấu hiểu: ``Dạ, cái này em hoàn toàn đoán được ạ. Vậy cụ thể thì công việc chính của em khi vào đây sẽ là gì thế, A-san?''

A-chan đẩy nhẹ gọng kính, nét mặt trở nên nghiêm túc hơn: ``Cậu còn nhớ tờ đơn đăng ký tuyển dụng mà cậu điền vào chiều nay chứ?''

``Vâng, em nhớ rất rõ. Có vấn đề gì sai sót hay thiếu thông tin gì hả chị?'' cậu nhíu mày lo lắng.

``Thực ra là\dots\ phía công ty chúng tôi có viết thiếu một số nội dung cực kỳ quan trọng trong mô tả công việc. Mà cái này cũng tại Natsuiro-san với Shirakami-san cứ liên tục giục giã, bắt tôi phải gửi tờ đơn đó lên web sớm quá nên mới xảy ra sơ suất như vậy\dots'' A-san liếc nhìn hai cô nàng kia đầy trách móc.

Kuon nhìn sang phía Matsuri đang giả vờ ngó lơ đi chỗ khác và Fubuki đang cười trừ trừ đầy vô tội, cậu nghiêng đầu khó hiểu: ``Ý chị là sao ạ? Thiếu nội dung gì thế chị?''

Cô từ tốn giải thích: ``Ngoài công việc hỗ trợ kỹ thuật IT thông thường ra, cậu sẽ còn phải đảm nhận một nhiệm vụ khác, có vai trò quan trọng không kém gì đâu.''

``Và công việc đó cụ thể là\dots?'' cậu càng hoang mang hơn.

Ngay lập tức, Matsuri hào hứng vung tay lên cao, hớn hở cướp lời: ``Đó là trở thành một phần cốt lõi của gia đình \hll\ tụi em, đồng hành và giúp đỡ tụi em thực hiện những ước mơ to lớn nhất!''

Kuon ngồi ngây người ra trước màn hình, chỉ biết phát ra một tiếng: ``Hả?'' đầy ngơ ngác và bất ngờ.

CEO Tanigo ôn tồn giải thích rõ ràng hơn: ``Hiện tại vị trí Tổng Quản Lý (General Manager) của chúng tôi đang bị bỏ trống, và cậu chính là ứng viên sáng giá được chọn cho vị trí đó đấy.''

Kuon tròn mắt ngạc nhiên, lắp bắp hỏi: ``Ờm\dots\ tại sao lại là cháu ạ? Cháu thấy mình đâu có kinh nghiệm gì đâu.''

A-chan thở dài giải thích: ``Vì chức vụ này chỉ tuyển đúng một chỉ tiêu duy nhất, và biểu mẫu đăng ký mà cậu điền vào chính là tờ đơn duy nhất mà chúng tôi đã thiết kế và thử nghiệm hệ thống.''

``Nghĩa là\dots\ kiểu ai nhanh tay đăng ký trước thì người đó được nhận luôn ạ?'' cậu con trai há hốc miệng kinh ngạc.

``Có thể hiểu nôm na là như vậy đấy,'' chị ấy nhún vai.

``Em không ngờ hệ thống lại hoạt động hiệu quả và nhanh đến thế luôn á!'' Matsuri phấn khích reo lên.

Fubuki vênh mặt lên đầy tự hào: ``Ý tưởng đăng đơn tuyển dụng thử nghiệm này là của tôi đề xuất mà lị! Quá là đỉnh cao luôn!''

A-chan liếc xéo hai cô nàng: ``Hai người không cần phải tranh nhau nhận công lao ở đây đâu nhé\dots''

Kuon sững sờ mất vài giây để cố gắng tiêu hóa và chấp nhận cái sự thật không tưởng vừa nghe. Đây chắc chắn là một cú sốc cực lớn trong cuộc đời cậu. Có lẽ trong cả ngành công nghiệp giải trí, chưa từng có một công ty nào lại đi tuyển dụng nhân sự cấp cao theo một phương thức kỳ quặc và `tùy duyên' đến mức độ này.

``Nghĩa là\dots\ cháu sẽ vừa phải làm kỹ thuật IT bảo trì máy móc, vừa kiêm luôn vai trò tổng quản lý chăm lo cho các idol luôn ạ!?'' cậu hỏi lại để chắc chắn mình không nghe nhầm.

``Về mặt lý thuyết thì đúng là như vậy đấy cậu Kuon,'' Tanigo gật đầu khẳng định.

``Vãi\dots'' Kuon há hốc miệng cạn lời, ``Nhưng cháu thực sự chưa từng có bất kỳ kinh nghiệm hay kiến thức nào về việc quản lý và vận hành idol cả!''

A-chan mỉm cười trấn an: ``Vấn đề đó cậu không cần lo lắng quá đâu. Tôi sẽ là người trực tiếp hướng dẫn và chỉ dạy dần cho cậu từ những bước cơ bản nhất. Mà tôi đoán cậu cũng đã có hình dung sơ bộ về công việc của một quản lý rồi đúng chứ?''

Cậu gãi đầu ngẫm nghĩ: ``Chắc là\dots\ ghi chép sổ sách, lập kế hoạch lịch trình phát sóng, giám sát hỗ trợ các idol trong các sự kiện, đại loại thế ạ?''

``Rất chính xác. Nhưng ngoài ra, cậu sẽ còn phải phụ trách cả việc bảo trì phần mềm, lắp đặt và kiểm tra hệ thống âm thanh ánh sáng sân khấu, chạy thử camera thu hình 3D, và sửa chữa các thiết bị máy móc chuyên dụng khác nữa. Tôi dám chắc một dân chuyên công nghệ như cậu sẽ rành mấy cái khoản này chứ?'' cô nàng Đeo kính nhướn mày dò hỏi.

Kuon hơi ngập ngừng đáp: ``Phần kỹ thuật này thì em không dám nhận là mình quá giỏi, nhưng mấy món lắp ráp, cài đặt phần mềm và sửa chữa cơ khí cơ bản thì em tự tin làm được ạ.''

``Thế là quá đủ rồi. Cứ bắt tay vào làm rồi mọi thứ sẽ quen dần thôi mà,'' A-chan gật đầu hài lòng.

Cậu khẽ nhíu mày suy nghĩ: ``Nghe sơ qua thế này là em biết khối lượng công việc sắp tới sẽ cực kỳ khổng lồ rồi đây\dots''

Roboco mỉm cười nhẹ nhàng qua màn hình: ``Em hơi bị ngạc nhiên là anh nghe xong đống nhiệm vụ đó mà không thấy phát hoảng hay bỏ chạy đấy nhé.''

Kuon nhún vai đầy tự tin: ``Anh chỉ sợ đúng hai thứ: một là lịch làm việc bị trùng với lịch học đại học của anh, hai là lương không đủ đóng học phí thôi. Còn lại thì\dots\ cứ giao hết đây anh xử lý tuốt!''

Nghe câu trả lời đầy khí thế của cậu, nàng Đeo kính nở một nụ cười rạng rỡ đầy an tâm: ``Phấn chấn và tràn đầy năng lượng thế là rất tốt đấy. Về phần lịch trình thì cậu cứ yên tâm, tôi sẽ chủ động sắp xếp linh hoạt và tạo điều kiện tối đa dựa theo lịch học của cậu ở trường. Cậu hoàn toàn có thể chọn làm bán thời gian hoặc toàn thời gian tùy ý.''

Nghe người đồng cấp tương lai nói vậy, Kuon thầm thở phào nhẹ nhõm trong lòng. Bất chợt, như nghĩ ra điều gì đó, cậu ngẩng đầu lên hỏi: ``À chị A này ơi, công ty mình có làm việc vào hai ngày thứ Bảy và Chủ Nhật không ạ?''

Aqua tròn mắt ngạc nhiên trước câu hỏi kỳ lạ của cậu, e dè lên tiếng: ``Ơ\dots\ Sao anh lại hỏi chuyện làm việc cuối tuần thế ạ?''

Cậu bình thản đáp: ``Tại vì anh muốn đăng ký làm tăng ca thêm vào hai ngày cuối tuần đó nữa.''

Câu trả lời dứt khoát của Kuon khiến cả phòng họp Zoom, kể cả vị CEO Tanigo, đều phải đứng hình mất vài giây và nhìn cậu bằng mắt đầy kinh ngạc. Dẫu biết rằng ở Nhật Bản, văn hóa làm việc tăng ca là chuyện cực kỳ phổ biến, thậm chí đôi khi còn là bắt buộc đối với nhân viên công sở, nhưng\dots

``\dots Tới mức chủ động xin làm thêm siêng năng như thế này thì quả thực chúng tôi chưa từng thấy bao giờ luôn đấy,'' A-chan đổ mồ hôi hột, dở khóc dở cười bình luận qua màn hình. ``Nhưng nếu làm liên tục như vậy, cậu không sợ bản thân sẽ bị kiệt sức và quá tải sao?''

Kuon lắc đầu cười đáp: ``Chị yên tâm, riêng hai ngày cuối tuần đó thì em xin phép được về sớm hơn một chút ạ. Em tự biết giới hạn sức khỏe của bản thân mình mà. Như thế có được chấp thuận không hả chị?''

``Cái đó cũng còn tùy thuộc vào khối lượng công việc cụ thể phát sinh nữa. Nhưng nếu cậu đã có nguyện vọng như vậy thì phía công ty hoàn toàn có thể tạo điều kiện tối đa. Tất nhiên, mọi giờ làm thêm ngoài giờ hành chính đều sẽ được tính lương tăng ca đầy đủ theo đúng luật lao động,'' A-san gật đầu đồng ý.

CEO Tanigo lúc này mới lên tiếng hỏi: ``Nhân tiện nhắc đến chuyện tiền lương, tôi muốn hỏi một chút. Chi phí sinh hoạt và học tập mỗi tháng của cậu khoảng bao nhiêu?''

Kuon nhanh chóng nhẩm tính rồi trả lời: ``Dạ, cháu chuẩn bị nhập học đại học ạ. Trường của cháu đóng học phí theo từng kỳ, mỗi kỳ kéo dài ba tháng và số tiền phải đóng ước chừng khoảng mười triệu đồng tiền Etsunan, quy đổi ra tiền Nhật thì xấp xỉ khoảng 50,000 yên ạ. Còn chi phí tiêu xài cá nhân hàng ngày của cháu thì ít lắm, cháu không có nhu cầu mua sắm gì nhiều đâu ạ.''

Tanigo nghe xong liền nở một nụ cười vô cùng rạng rỡ, ông đập tay xuống bàn đầy hào hứng: ``Ồ! Thế thì cậu cứ yên tâm tuyệt đối nhé! Với mức lương cơ bản của chức vụ Tổng Quản Lý này, chưa đầy hai tháng làm việc của cậu là đã thừa sức chi trả toàn bộ học phí cho cả một học kỳ rồi. Đó là còn chưa tính đến khoản tiền thưởng tăng ca cuối tuần nữa đấy, tính ra là dư dả vô cùng luôn!''

Kuon nghe tới đó thì hai mắt sáng lên đầy phấn khích: ``Cháu nghe thấy chế độ đãi ngộ có vẻ cực kỳ hậu hĩnh và hấp dẫn đấy nhé! Chú Tanigo nói là phải giữ lời hứa đấy nha!''

``Tôi là CEO cơ mà, danh dự của tôi đảm bảo sẽ không bao giờ phản bội hay đối xử bất công với nhân viên của mình đâu mà cậu phải sợ. Chỉ cần cậu làm việc thật tốt và có trách nhiệm, lương thưởng chắc chắn sẽ không khiến cậu phải thất vọng đâu. Được chứ?'' Tanigo cười lớn đầy sảng khoái.

Kuon tràn đầy khí thế hứa hẹn: ``Dạ, cháu xin hứa sẽ cống hiến hết mình ạ!''

``Vậy là các điều khoản cơ bản về lịch trình làm việc và tiền lương cũng đã thống nhất xong xuôi rồi nhé. Sáng ngày mai là Chủ Nhật đúng không nhỉ? Cậu có bận gì không?'' Tanigo hỏi.

``Dạ không ạ, hôm đó cháu hoàn toàn rảnh rỗi ạ.''

``Tốt quá. Vậy sáng mai cậu cứ đến trực tiếp văn phòng công ty theo địa chỉ tôi gửi. Chúng ta sẽ gặp mặt trao đổi trực tiếp tại văn phòng CEO để chốt lại hợp đồng và ký kết chính thức luôn. OK chứ?''

``Dạ vâng ạ, cháu nhất định sẽ đến đúng giờ!'' Kuon lễ phép đáp.

A-chan mỉm cười nhẹ nhõm: ``Tuyệt vời quá, vậy là từ ngày mai tôi sẽ chính thức có thêm một người đồng nghiệp đáng tin cậy để chia sẻ bớt gánh nặng công việc rồi.''

``Dạ vâng, sau này có gì sơ suất kính mong chị chỉ giáo và giúp đỡ em nhiều hơn nhé, A-san!'' cậu cúi đầu chào lịch sự qua camera.

Cả bốn cô gái Roboco, Fubuki, Aqua và Matsuri cũng đồng loạt vẫy tay chào tạm biệt đầy nhí nhảnh qua màn hình: ``Nhờ anh giúp đỡ và hướng dẫn tụi em trong thời gian tới nhé, anh quản lý mới!''

``Cảm ơn mọi người nhiều nhé. Anh nhất định sẽ cố gắng hết sức mình!'' cậu mỉm cười vẫy tay chào lại, lòng tràn đầy quyết tâm.

\begin{center}
  {\textit{*Các nhân vật nói bằng tiếng Etsunan.}}
\end{center}

Ngay sau khi cuộc gọi kết thúc, Kuon quay sang bố, người vừa nghe hết những gì mà cậu và mọi người đã bàn bạc. ``Bố à, con vừa nhận được công việc rồi! Chỉ còn ký hợp đồng nữa thôi là xong. Lương cũng cao lắm!'' cậu hào hứng nói. Người bố gật đầu, khẽ mỉm cười tự hào: ``Thế là từ nay con tự lo được cho mình rồi đấy nhỉ?''

Không bỏ lỡ cơ hội, ông kéo ra vài chai bia cất ở trong hộc tủ, vỗ vai con trai: ``Nào, khui bia ăn mừng! Đây là bước ngoặt của đời con đấy.'' Hai bố con cùng ngồi xuống bàn, nâng cốc chúc mừng cho khởi đầu mới. Kuon nở nụ cười nhẹ, cảm nhận sự ấm áp từ ánh mắt tự hào của bố - người bạn vong niên của cậu.

Đêm đó, sau ``đại hội bia'' kéo dài tới tận khuya, Kuon nằm trằn trọc trên giường. Dù hơi men đã tan, tâm trí cậu vẫn không ngừng nghĩ về những điều sắp tới. Hình ảnh văn phòng làm việc, các cô idol và cả ông CEO kỳ lạ hiện lên như một giấc mơ lạ lùng.

``Ngày mai\dots\ mình sẽ làm được,'' cậu khẽ thì thầm.

Ngày hôm sau, cậu đến văn phòng theo chỉ dẫn của ông CEO để bắt đầu nhận việc, mà không hề biết rằng\dots

\dots cậu sẽ trở thành một ``nam chính'' trong một vở kịch đầy hài hước và hỗn loạn của công ty giải trí ảo vươn tầm thế giới mang tên: \textbf{Holo no Graffiti.}

Và giờ đây, màn kịch của chúng ta xin phép được\dots

\begin{center}
  \textbf{\huge BẮT ĐẦU.}
\end{center}
\end{document}